% Auto-generated by Code2Latex
% Requires: longtable, booktabs packages

\begin{longtable}{p{\textwidth}}
\caption{Tools for single (Level 1)} \\
\endfirsthead

\multicolumn{1}{c}{\tablename\ \thetable{} -- continued from previous page} \\
\endhead

\multicolumn{1}{r}{Continued on next page} \\
\endfoot

\endlastfoot

\textbf{\texttt{calculate\_parallel\_resistance}} \\
\begin{description}
    \item[\textbf{Description:}] Calculate total resistance of resistors connected in parallel.

Computes the equivalent resistance when multiple resistors are connected
across the same two nodes. In parallel configuration, current divides among the
resistors, and the reciprocal of total resistance equals the sum of reciprocals
of individual resistances. This calculation is crucial for simplifying parallel branches in a circuit.
    \item[\textbf{Arguments:}]
    \begin{description}
        \item \texttt{resistances} (list[float]): List of resistance values in ohms.
A list containing floating-point numbers, each representing the resistance of an individual resistor. All values must be positive.
    \end{description}
    \item[\textbf{Returns:}] Total parallel resistance in ohms.

A single floating-point number representing the equivalent resistance of all input resistors connected in parallel.
\end{description} \\
\midrule
\textbf{\texttt{calculate\_series\_resistance}} \\
\begin{description}
    \item[\textbf{Description:}] Calculate total resistance of resistors connected in series.

Computes the equivalent resistance when multiple resistors are connected
end-to-end in a single path. In series configuration, current flows through each
resistor sequentially, and the total resistance is the sum of individual resistances.
    \item[\textbf{Arguments:}]
    \begin{description}
        \item \texttt{resistances} (list[float]): List of resistance values in ohms.
A list containing floating-point numbers, each representing the resistance of an individual resistor. All values must be positive.
    \end{description}
    \item[\textbf{Returns:}] Total series resistance in ohms.

A single floating-point number representing the sum of all input resistances.
\end{description} \\
\midrule
\textbf{\texttt{cat\_files}} \\
\begin{description}
    \item[\textbf{Description:}] Concatenate and display contents of one or more files
    \item[\textbf{Arguments:}]
    \begin{description}
        \item \texttt{paths} (list[str]): List of file paths to concatenate
        \item \texttt{separator} (str): Separator between file contents
    \end{description}
\end{description} \\
\midrule
\textbf{\texttt{copy\_file}} \\
\begin{description}
    \item[\textbf{Description:}] Copy a file from a source to destination
    \item[\textbf{Arguments:}]
    \begin{description}
        \item \texttt{source} (str): Source file path
        \item \texttt{destination} (str): Destination file path
    \end{description}
\end{description} \\
\midrule
\textbf{\texttt{delta\_to\_wye\_transform}} \\
\begin{description}
    \item[\textbf{Description:}] Convert delta (triangle) resistor configuration to wye (star) configuration.

Transforms a three-resistor delta network into an equivalent three-resistor
wye network. This is essential for solving complex resistor networks that cannot be
reduced using simple series/parallel combinations. The transformation preserves the
resistance between any two external nodes, simplifying nodal analysis.
    \item[\textbf{Arguments:}]
    \begin{description}
        \item \texttt{ra} (float): Resistance between nodes A and B in delta configuration.
A positive floating-point number representing the resistance of the resistor connected between nodes A and B in the delta network.
        \item \texttt{rb} (float): Resistance between nodes B and C in delta configuration.
A positive floating-point number representing the resistance of the resistor connected between nodes B and C in the delta network.
        \item \texttt{rc} (float): Resistance between nodes C and A in delta configuration.
A positive floating-point number representing the resistance of the resistor connected between nodes C and A in the delta network.
    \end{description}
    \item[\textbf{Returns:}] JSON string with keys 'r1', 'r2', 'r3' for wye resistor values.

A JSON string containing a dictionary with three keys: 'r1', 'r2', and 'r3', whose values are the calculated equivalent resistances for the wye network, each being a floating-point number. R1 is connected to original node A, R2 to B, and R3 to C.
\end{description} \\
\midrule
\textbf{\texttt{file\_info}} \\
\begin{description}
    \item[\textbf{Description:}] Get information about a file or directory
    \item[\textbf{Arguments:}]
    \begin{description}
        \item \texttt{path} (str): Path to the file or directory
    \end{description}
\end{description} \\
\midrule
\textbf{\texttt{list\_files}} \\
\begin{description}
    \item[\textbf{Description:}] List files in a directory
    \item[\textbf{Arguments:}]
    \begin{description}
        \item \texttt{path} (str): Path to the directory
        \item \texttt{recursive} (bool): Whether to list files recursively
    \end{description}
\end{description} \\
\midrule
\textbf{\texttt{read\_file}} \\
\begin{description}
    \item[\textbf{Description:}] Read the contents of a file into a string
    \item[\textbf{Arguments:}]
    \begin{description}
        \item \texttt{path} (str): Path to the file to read
    \end{description}
\end{description} \\
\midrule
\textbf{\texttt{simulate\_circuit\_resistance}} \\
\begin{description}
    \item[\textbf{Description:}] Simulate total resistance between specified terminals in a given circuit topology.

Takes a circuit topology description and calculates the equivalent resistance
between two specified terminal nodes using nodal analysis. This tool allows for the
testing and validation of topology hypotheses by comparing calculated resistance with
measured values. It provides a foundational capability for circuit analysis and design.
    \item[\textbf{Arguments:}]
    \begin{description}
        \item \texttt{topology} (str): JSON string describing circuit.
A JSON string that defines the circuit's components and their interconnections. It must contain a "resistors" dictionary (mapping resistor IDs to their resistance values) and a "connections" list (each entry being a list `[node1, node2, resistor\_id]`).
        \item \texttt{terminal\_nodes} (list[str]): List of two node names to measure resistance between.
A list containing exactly two strings, where each string is the name of a node in the circuit. The tool will calculate the equivalent resistance between these two specified nodes.
    \end{description}
    \item[\textbf{Returns:}] Equivalent resistance between terminals in ohms.

A floating-point number representing the total equivalent resistance measured between the two specified terminal nodes in the given circuit topology.
\end{description} \\
\midrule
\textbf{\texttt{validate\_measurements}} \\
\begin{description}
    \item[\textbf{Description:}] Validate if proposed topology matches all given measurements.

Compares the resistance predictions of a proposed circuit topology against
actual measurements. It quantifies the discrepancy by calculating error metrics,
which are essential for validating hypotheses about circuit structure and resistor values.
This tool helps in refining and confirming circuit designs.
    \item[\textbf{Arguments:}]
    \begin{description}
        \item \texttt{topology} (str): JSON string describing proposed circuit topology.
A JSON string conforming to the `CircuitTopology` structure, including resistor IDs, their estimated values, and the connections between nodes. This represents your hypothesis about the circuit's structure.
        \item \texttt{measurements} (str): JSON string with actual measurements.
A JSON string representing a list of CircuitMeasurement objects. Each object should contain node\_a, node\_b, and at least resistance. These are the real-world observations.
    \end{description}
    \item[\textbf{Returns:}] JSON string with validation results including error metrics.

A JSON string containing a dictionary with various error metrics: `total\_error`, `max\_error`, `mean\_error`, and `detailed\_errors` (a list of per-measurement errors including predicted, actual, absolute error, and relative error). It also includes `num\_measurements`. This output helps quantify the accuracy of the proposed topology.
\end{description} \\
\midrule
\textbf{\texttt{write\_file}} \\
\begin{description}
    \item[\textbf{Description:}] Write content to a file
    \item[\textbf{Arguments:}]
    \begin{description}
        \item \texttt{path} (str): Path to the file to write
        \item \texttt{content} (str): Content to write to the file
    \end{description}
\end{description} \\
\midrule
\textbf{\texttt{wye\_to\_delta\_transform}} \\
\begin{description}
    \item[\textbf{Description:}] Convert wye (star) resistor configuration to delta (triangle) configuration.

Transforms a three-resistor wye network into an equivalent three-resistor
delta network. This is the inverse of delta-to-wye transformation and is useful when
the delta form provides easier analysis or when testing different topology hypotheses. This transformation is key for circuit simplification in cases where a wye configuration is encountered.
    \item[\textbf{Arguments:}]
    \begin{description}
        \item \texttt{r1} (float): Wye resistor connected to node A.
A positive floating-point number representing the resistance of the resistor connected from the center of the wye to node A.
        \item \texttt{r2} (float): Wye resistor connected to node B.
A positive floating-point number representing the resistance of the resistor connected from the center of the wye to node B.
        \item \texttt{r3} (float): Wye resistor connected to node C.
A positive floating-point number representing the resistance of the resistor connected from the center of the wye to node C.
    \end{description}
    \item[\textbf{Returns:}] JSON string with keys 'ra', 'rb', 'rc' for delta resistor values.

A JSON string containing a dictionary with three keys: 'ra', 'rb', and 'rc', whose values are the calculated equivalent resistances for the delta network, each being a floating-point number. Ra is between original nodes A and B, Rb between B and C, and Rc between C and A.
\end{description} \\
\end{longtable}
